\RequirePackage{pdfmanagement}
\SetKeys[document/metadata]{pdfversion=1.7, pdfstandard=A-2u, lang=en-GB}
\input{preamble/pdfA-compliance/glyphtounicode}
\documentclass[11pt,a4paper]{report}

\newcommand*{\PaperTitle}{Taylor dispersion in regular polygonal ducts: the sign of the skewness at all times}
\newcommand*{\PaperTitlePlain}{Taylor dispersion in regular polygonal ducts: the sign of the skewness at all times}
\newcommand*{\PaperShortTitle}{Skewness in regular polygonal ducts at all times}
\newcommand*{\PaperSubtitle}{chemistry-open-problems \textperiodcentered{} 31 October 2026 \textperiodcentered{} preliminary, not yet verified}
\newcommand*{\PaperCredits}{Solver: Claude Opus 5.5 (effort not recorded) \textperiodcentered{} Independent verification: pending}
\newcommand*{\PaperKeywords}{Taylor dispersion, skewness, regular polygon, Neumann heat semigroup, torsion function, interval arithmetic}
\input{chemop-preamble}
\addbibresource{references.bib}

\DeclareRobustCommand*{\avg}[1]{\langle #1\rangle}
\DeclareRobustCommand*{\norm}[1]{\lVert #1\rVert}
\DeclareRobustCommand*{\bdnorm}[1]{\lVert #1\rVert_{\partial}}
\DeclareRobustCommand*{\dualnorm}[1]{\lVert #1\rVert_{*}}
\DeclareRobustCommand*{\ut}{\tilde u}
\DeclareRobustCommand*{\dd}{\mathop{}\!\mathrm{d}}
\DeclareRobustCommand*{\Pg}{P_n}
\DeclareRobustCommand*{\sg}{\mathrm{e}^{s\Delta}}

\begin{document}
\makepapertitle

\begin{paperabstract}
Aminian, Camassa and McLaughlin claimed from finite-element computations that among regular polygonal ducts the
equilateral triangle is the only one whose solute skewness changes sign, ``all other polygons possess positive
skewness at all times'', adding that this ``cannot be proved on finite times''. We prove the claim. For every regular
\(n\)-gon with \(n\ge4\), every Péclet number and every time, the skewness of a cross-sectionally uniform release with
symmetric longitudinal profile is strictly positive; the same holds for the disk, and for the equilateral triangle
the skewness changes sign exactly once. The proof has four parts. An identity writes the third moment as
\(K(\tau)=2\int_0^\tau\langle u\,A(s)A(\tau-s)\rangle\dd s\), where \(u\) is the mean-free Poiseuille profile and \(A\)
is the time integral of the Neumann heat flow of \(u\). A maximum-principle argument, valid on every convex section,
gives an explicit lower bound for \(K\) that is positive on \(0<\tau\le0.015\) for all \(n\) at once, even though
\(\langle u^3\rangle\to0\) as \(n\to\infty\). A modal approximation of \(A\) with an a-posteriori energy estimate certifies
\(K>0\) on \([0.015,\infty)\) in ball arithmetic for \(4\le n\le199\). For \(n\ge200\) the same estimate is applied
with Bessel modes and closed-form disk data, with errors bounded uniformly in \(n\). The paper also collects the
long- and short-time results of the same problem.
\end{paperabstract}

\section*{The question}

M.~Aminian, R.~Camassa and R.~M.~McLaughlin~\autocite{ACM2018,ACM2017} study a passive tracer released uniformly across
a straight duct in pressure-driven flow. The sign of the skewness of its cross-sectionally averaged distribution
says whether the tracer is front- or back-loaded. The abstract of the arXiv version (v1--v3) says:
\begin{openquote}
``Interestingly, the equilateral triangle is the only regular polygon with this property, all other polygons possess
positive skewness at all times, although this cannot cannot be proved on finite times due to the lack of closed form
flow solutions for such geometries.''
\end{openquote}
``This property'' means positive short-time and negative long-time skewness. The published abstract weakens the
sentence to ``The equilateral triangle appears to be the only regular polygon with this property---all others possess
positive skewness at all times.'' The evidence in the source is finite-element computation of the two asymptotic
coefficients for finitely many \(n\). An earlier stage of the present work~\autocite{Result14} settled the two
asymptotic signs (\zcref{thm:ends}). The sign at intermediate times was open.

Let \(\Omega\) be the section, \(\ut\) the Poiseuille profile (\(-\Delta\ut=2\), \(\ut=0\) on \(\partial\Omega\)),
\(u=\ut-\avg{\ut}\), and \(P_s=\sg\) the Neumann heat semigroup on \(\Omega\). Put
\begin{equation}\label{eq:K}
A(s)=\int_0^sP_ru\dd r,\qquad K(\tau)=2\int_0^\tau\avg{u\,A(s)\,A(\tau-s)}\dd s .
\end{equation}
By \zcref{prop:moments}, the third central moment of the release is \(\mu_3(\tau)=\kappa_3+3\mathrm{Pe}^3K(\tau)\),
where \(\kappa_3\) is the third cumulant of the initial longitudinal profile, and \(\mu_2(\tau)>0\). So for a symmetric
longitudinal profile, in particular a plane source, the sign of the skewness is the sign of \(K(\tau)\) for every
Péclet number.

\begin{theorem}\label{thm:main}
For every integer \(n\ge4\) and the regular \(n\)-gon, \(K(\tau)>0\) for all \(\tau>0\). Hence the skewness is strictly
positive at all times, for every Péclet number and every cross-sectionally uniform release with symmetric longitudinal
profile. The same holds for the disk.
\end{theorem}

\begin{theorem}\label{thm:triangle}
For the equilateral triangle \(K\) has exactly one zero \(\tau^*\) on \((0,\infty)\), simple, with \(K>0\) before and
\(K<0\) after; \(\tau^*\in[1.17,1.25]\) in the time unit \(a^2/\kappa\), \(a\) the inradius.
\end{theorem}

\begin{theorem}[note=the two ends]\label{thm:ends}
Let \(g\) be the mean-free solution of \(-\Delta g=u\), \(\partial_ng=0\). Then
\(K(\tau)=2\avg{ug^2}\tau+O(1)\) as \(\tau\to\infty\) and \(K(\tau)=\avg{u^3}\tau^3/3+O(\tau^{7/2})\) as \(\tau\to0\), and
\begin{enumerate}
\item \(\avg{ug^2}=-7101/476476000<0\) for \(n=3\), and \(\avg{ug^2}>0\) for every \(n\ge4\) and for the disk;
\item \(\avg{u^3}=27/19250\) for \(n=3\) (inradius 1), \(\avg{u^3}>0\) for every \(n\ge4\), \(\avg{u^3}=0\) for the disk, and
\(\avg{u^3}=\tan^4(\pi/n)/360+O(\tan^5(\pi/n))\).
\end{enumerate}
\end{theorem}

In the notation of the source \(2\avg{ug^2}=\avg{ug_2}\), see~\eqref{eq:Sinf}. The certified quantities behind \zcref{thm:main} reprove item~1 of \zcref{thm:ends}
independently of~\autocite{Result14}. \zcref{sec:table} gives enclosures of all entries of Table~1 of the source.

\begin{remark}
\zcref{thm:main} needs \(\kappa_3=0\). For a non-symmetric longitudinal profile
\(\mu_3=\kappa_3+3\mathrm{Pe}^3K\), and the sign depends on the data. The results concern the moments of the
Aris hierarchy, exactly as in the source.
\end{remark}

\section{Exact formulas}\label{sec:formulas}

Throughout, \(\Pg\) is the regular \(n\)-gon with apothem \(1\) and side \(\{x=1,\ |y|\le t\}\), \(t=\tan(\pi/n)\),
\(R=\sec(\pi/n)=\sqrt{1+t^2}\), \(|\Pg|=nt\), \(|\partial\Pg|=2nt\). For a section \(\Omega\) we write
\(\avg f=|\Omega|^{-1}\int_\Omega f\), \(\norm f^2=\avg{f^2}\), \(\bdnorm f^2=|\Omega|^{-1}\oint_{\partial\Omega}f^2\).
Further \(c=\avg{\ut}\), \(u=\ut-c\), \(m=-\partial_n\ut\ge0\) (with \(|\Omega|^{-1}\oint m=2\)). All statements are
invariant under dilations, with \(\tau\) scaling like a length squared. We write \(g\) for the function \(g_1\) of the
source and \(\mathcal T(f,h)=\avg{ufh}\).

The tracer solves \(T_\tau+\mathrm{Pe}\,uT_x=\Delta_3T\) in \(\mathbb R\times\Omega\), \(\partial_nT=0\),
\(T(x,\cdot,0)=\varrho(x)\), \(\varrho\ge0\), \(\int\varrho=1\), with finite moments, and
\(\int x\varrho=0\) after a translation. The moments \(c_p(\cdot,\tau)=\int x^pT\dd x\) solve the Aris hierarchy
\(\partial_\tau c_p=\Delta c_p+p(p-1)c_{p-2}+p\,\mathrm{Pe}\,uc_{p-1}\), \(\partial_nc_p=0\), \(c_p(0)=\int x^p\varrho\)~\autocite{Aris1956}.

\begin{proposition}\label{prop:moments}
\(\mu_2(\tau)>0\) and \(\mu_3(\tau)=\kappa_3+3\mathrm{Pe}^3K(\tau)\) with \(K\) as in~\eqref{eq:K}. Moreover, with
\(G(a,b)=\avg{u\,P_au\,P_bu}\),
\begin{equation}\label{eq:KG}
K(\tau)=2\iint_{a,b\ge0,\ a+b\le\tau}(\tau-a-b)\,G(a,b)\dd a\dd b ,
\end{equation}
\(A(s)=g-P_sg\), and \(K(\tau)=2\avg{ug^2}\tau+O(1)\) as \(\tau\to\infty\). If \(g_2\) is the mean-free solution of
\(-\Delta g_2=2(ug-\avg{ug})\), \(\partial_ng_2=0\), then \(\avg{ug_2}=2\avg{ug^2}\) and
\begin{equation}\label{eq:Sinf}
\tau^{1/2}\mathrm{Sk}(\tau)\to\frac{3\mathrm{Pe}^3\avg{ug_2}}{(2+2\mathrm{Pe}^2\avg{ug})^{3/2}}\ ,
\end{equation}
which is equation~(17) of the source.
\end{proposition}
\begin{proof}
\(c_0=1\), \(c_1=\mathrm{Pe}\,A\), \(\avg{c_1}=0\). Next \(c_2=\int x^2\varrho+2\tau+\mathrm{Pe}^2S\) with
\(S'=\Delta S+2uA\), \(S(0)=0\). Hence \(\mu_2'=2+2\mathrm{Pe}^2\avg{uA}\) with
\(\avg{uA(s)}=\int_0^s\norm{P_{r/2}u}^2\dd r\ge0\), so \(\mu_2>0\). Further \(\frac{\dd}{\dd\tau}\avg{c_3}=6\avg{c_1}+3\mathrm{Pe}\avg{uc_2}
=3\mathrm{Pe}^3\avg{uS}\). By Duhamel and self-adjointness, \(\avg{uS(\sigma)}=2\int_0^\sigma\avg{A'(\sigma-s)\,u\,A(s)}\dd s\).
Integrating over \(\sigma\in[s,\tau]\) with \(A(0)=0\) gives~\eqref{eq:K}, and~\eqref{eq:KG} follows by inserting
\(A(s)=\int_0^sP_ru\dd r\): for \(r+q\le\tau\) the variable \(s\) runs over an interval of length \(\tau-r-q\).
Spectrally \(A(s)=g-P_sg\). Insert this into~\eqref{eq:K}. The terms containing \(P_sg\) stay bounded because of the
spectral gap: \(\int_0^\tau|\mathcal T(g,P_sg)|\dd s\) is bounded, and \(|\mathcal T(P_sg,P_{\tau-s}g)|\le C\mathrm e^{-\nu\tau}\).
So \(K=2\avg{ug^2}\tau+O(1)\). For small \(\tau\), \(\norm{P_au-u}\le0.64\sqrt a\norm{\nabla u}\) spectrally and \(u\) is bounded, so \(|G(a,b)-\avg{u^3}|\le C(\sqrt a+\sqrt b)\) and~\eqref{eq:KG} gives \(K=\avg{u^3}\tau^3/3+O(\tau^{7/2})\). Testing the two weak forms against \(g_2\) and \(g\) gives
\(\avg{ug_2}=\avg{\nabla g\cdot\nabla g_2}=2\avg{ug^2}\); and \(\mu_2\sim(2+2\mathrm{Pe}^2\avg{ug})\tau\).
\end{proof}

\section{Tools}\label{sec:tools}

Let \(V\) be the space of \(D_n\)-invariant, mean-free functions in \(H^1(\Pg)\). Such a function is determined by its
restriction to the half-sector \(T=\{0\le y\le tx,\ x\le1\}\), a right triangle of diameter \(R\), and its restriction
is mean-free on \(T\). The functions \(u\), \(g\), \(P_su\), \(P_sg\), \(A(s)\) lie in \(V\). Payne--Weinberger on the convex set
\(T\)~\autocite{PayneWeinberger1960,Bebendorf2003} gives
\begin{equation}\label{eq:PW}
\norm\varphi\le C_P\norm{\nabla\varphi},\quad C_P=\frac R\pi,\qquad \norm{P_s\varphi}\le\mathrm e^{-\nu s}\norm\varphi,\quad
\nu=\frac{\pi^2}{R^2}\qquad(\varphi\in V).
\end{equation}
Since \(x\cdot n=1\) on \(\partial\Pg\), \(\bdnorm\varphi^2=\avg{\operatorname{div}(x\varphi^2)}=2\norm\varphi^2+2\avg{\varphi\,x\cdot\nabla\varphi}\), so
\begin{equation}\label{eq:trace}
\bdnorm\varphi\le\tau_{\mathrm{tr}}\norm{\nabla\varphi},\qquad\tau_{\mathrm{tr}}^2=2C_P^2+2RC_P .
\end{equation}
For \(R_*\in L^2\), \(b\in L^2(\partial\Pg)\) the functional \(\ell(\varphi)=\avg{R_*\varphi}+|\Pg|^{-1}\oint b\varphi\) on
\(V\) has dual norm \(\dualnorm\ell=\sup_{\varphi\in V}|\ell(\varphi)|/\norm{\nabla\varphi}\le C_P\norm{R_*-\avg{R_*}}+\tau_{\mathrm{tr}}\bdnorm b\).

\begin{lemma}\label{lem:harmonic}
Let \(E\in H^2(\Pg)\) be harmonic. Then \(\norm E^2\le\frac{1+R}2\bdnorm E^2\) and
\(\bdnorm{\partial_nE}\le(t+R)\bdnorm{\partial_sE}\).
\end{lemma}
\begin{proof}
\(\ut\in H^2(\Pg)\) by convexity~\autocite{Grisvard1985}. Green's identities give
\(2\int E^2=-2\int\ut|\nabla E|^2+\oint E^2m\). On a side, \(\ut\le(1-x)(x+R)\) by the maximum principle, so \(m\le1+R\).
The second bound is the Rellich identity: the field \((x\cdot\nabla E)\nabla E-\frac12|\nabla E|^2x\) is divergence-free. With
\(x\cdot n=1\) and \(x\cdot\nabla E=\partial_nE+\sigma\partial_sE\), \(|\sigma|\le t\), on a side, this gives
\(X^2\le Y^2+2tXY\) for \(X=\bdnorm{\partial_nE}\), \(Y=\bdnorm{\partial_sE}\).
\end{proof}

\begin{lemma}[note=cell problem]\label{lem:cell}
Let \(\ut_H\) be a \(D_n\)-invariant polynomial with \(-\Delta\ut_H=2\), and \(g_H\) a \(D_n\)-invariant polynomial with
\(-\Delta g_H=\ut_H-\avg{\ut_H}\). Then
\[
\norm{g-(g_H-\avg{g_H})}\le C_P\Bigl(C_P\bigl(\tfrac{1+R}2\bigr)^{1/2}\bdnorm{\ut_H}+\tau_{\mathrm{tr}}\bdnorm{\partial_ng_H}\Bigr).
\]
\end{lemma}
\begin{proof}
\(E=\ut-\ut_H\) is harmonic, in \(H^2\), and equals \(-\ut_H\) on \(\partial\Pg\). For \(w=g-g_H+\avg{g_H}\in V\) and \(\varphi\in V\),
\(\avg{\nabla w\cdot\nabla\varphi}=\avg{(E-\avg E)\varphi}-|\Pg|^{-1}\oint\varphi\,\partial_ng_H\). Use~\eqref{eq:PW},~\eqref{eq:trace} and
\zcref{lem:harmonic}.
\end{proof}

\begin{lemma}[note=modal approximation of \(A\)]\label{lem:modal}
Let \(\hat g\in V\), \(\psi_1,\dots,\psi_M\in V\), \(\mu_k>0\), \(d_k\in\mathbb R\), and let
\(\ell_k(\varphi)=\avg{\nabla\psi_k\cdot\nabla\varphi}-\mu_k\avg{\psi_k\varphi}\) on \(V\). Put
\(\tilde A(s)=\hat g-\sum_kd_k\mathrm e^{-\mu_ks}\psi_k\). If \(\norm{g-\hat g}\le\omega\), \(\norm{\hat g-\sum_kd_k\psi_k}\le\rho\) and
\(L_1=\sum_k|d_k|\dualnorm{\ell_k}(2\mu_k)^{-1/2}\), then
\[
\norm{A(s)-\tilde A(s)}\le\Delta_0:=2\omega+\rho+L_1\qquad(s\ge0).
\]
\end{lemma}
\begin{proof}
\(y(s)=P_sg-\sum_kd_k\mathrm e^{-\mu_ks}\psi_k\) satisfies \(\avg{y'\varphi}+\avg{\nabla y\cdot\nabla\varphi}=-\sum_kd_k\mathrm e^{-\mu_ks}\ell_k(\varphi)\)
on \(V\). With \(\varphi=y\), \(\frac{\dd}{\dd s}\norm y^2\le-2\norm{\nabla y}^2+2\ell(s)\norm{\nabla y}\le\ell(s)^2\), where
\(\ell(s)=\sum_k|d_k|\mathrm e^{-\mu_ks}\dualnorm{\ell_k}\). By Minkowski's inequality \(\int_0^\infty\ell^2\le L_1^2\), so
\(\norm{y(s)}\le\norm{y(0)}+L_1\le\omega+\rho+L_1\). Finally \(A(s)-\tilde A(s)=(g-\hat g)-y(s)\).
\end{proof}

\begin{lemma}[note=comparison]\label{lem:compare}
In the setting of \zcref{lem:modal} let \(w\) be bounded with \(\norm w_\infty\le U\), \(\norm{u-w}_\infty\le o\), and
\(\tilde K(\tau)=2\int_0^\tau\avg{w\tilde A(s)\tilde A(\tau-s)}\dd s\). Then \(|K(\tau)-\tilde K(\tau)|\le\mathcal E(\tau)\), where
\[
\mathcal E(\tau)=2U\bigl(2\Delta_0\min(\tfrac12\tau^2\norm u,\tau\norm g)+\Delta_0^2\tau\bigr)+2o\min(\tfrac16\tau^3\norm u^2,\tau\norm g^2).
\]
\end{lemma}
\begin{proof}
\(\norm{A(s)}\le\min(s\norm u,\norm g)\) by the spectral theorem. With \(e=A-\tilde A\),
\(\avg{uA(s)A(s')}-\avg{w\tilde A(s)\tilde A(s')}=\avg{(u-w)A(s)A(s')}+\avg{we(s)A(s')}+\avg{w\tilde A(s)e(s')}\) and
\(\norm{\tilde A}\le\norm A+\Delta_0\). Integrate.
\end{proof}

If \(\tilde A(s)=\sum_{k=0}^MD_k\mathrm e^{-\mu_ks}\psi_k\) with \(\mu_0=0\), \(\psi_0=\hat g\), \(D_0=1\), \(D_k=-d_k\), and
\(W_{kl}=D_kD_l\avg{w\psi_k\psi_l}\), then
\begin{equation}\label{eq:Ktilde}
\tilde K(\tau)=2\sum_{k,l=0}^MW_{kl}E_{kl}(\tau),\qquad E_{kl}(\tau)=\int_0^\tau\mathrm e^{-\mu_ka-\mu_l(\tau-a)}\dd a .
\end{equation}
Writing \(E_{kl}=\tau\int_0^1\mathrm e^{-\tau(\theta\mu_k+(1-\theta)\mu_l)}\dd\theta\) and using
\(\sup_{x\ge0}(3x^2+x^3)\mathrm e^{-x}<3\) gives \(|\partial_\tau^3E_{kl}|\le3\tau^{-2}\). So
\begin{equation}\label{eq:Taylor}
|\tilde K'''(\tau)|\le6\tau^{-2}\sum_{k,l}|W_{kl}|,
\end{equation}
and \(0\le E_{kl}\le1/\max(\mu_k,\mu_l)\) for \((k,l)\ne(0,0)\), so that \(\tilde K(\tau)\ge2W_{00}\tau-C_1\) with
\(C_1=2\sum_{(k,l)\ne(0,0)}|W_{kl}|/\max(\mu_k,\mu_l)\).

\section{Short times, uniformly in the section}\label{sec:short}

In this section \(\Omega\) is any bounded convex polygon or the disk. Let \(F(s)=\avg{uP_su}\). This is the only place
where the boundary layer is treated, and no discretisation enters.

\begin{lemma}\label{lem:layer}
\(\Gamma_s:=P_su-u+2s\ge0\) and \(\avg{\Gamma_s}=2s\). Moreover \(\Lambda(s):=\avg{\ut\,\Gamma_s}\) satisfies
\[
F(s)=\avg{u^2}-2cs+\Lambda(s),\qquad \Lambda(s)=\int_0^s|\Omega|^{-1}\oint_{\partial\Omega}m\,P_r\ut\dd\sigma\dd r .
\]
\end{lemma}
\begin{proof}
For \(\varphi\in H^1\) Green's identity gives \(\avg{\nabla u\cdot\nabla\varphi}=2\avg\varphi-|\Omega|^{-1}\oint m\varphi\). So
\(\avg{\partial_s\Gamma_s\,\varphi}+\avg{\nabla\Gamma_s\cdot\nabla\varphi}=|\Omega|^{-1}\oint m\varphi\), \(\Gamma_0=0\). Since \(m\ge0\),
testing with the negative part gives \(\Gamma_s\ge0\); testing with \(\varphi=1\) gives the mean. Next
\(\avg{\ut\Gamma_s}=\avg{uP_su}+c\avg{P_su}-\avg{\ut u}+2sc=F(s)-\avg{u^2}+2cs\). Finally
\(F'(r)=-\avg{\nabla u\cdot\nabla P_ru}=|\Omega|^{-1}\oint mP_ru=|\Omega|^{-1}\oint mP_r\ut-2c\).
\end{proof}

\begin{lemma}[note=barrier]\label{lem:barrier}
Let \(d\) be the distance to \(\partial\Omega\) and suppose \(\ut\le\bar m\,d+\beta\) on \(\Omega\). Then
\(\Lambda(s)\le\frac{8}{3\sqrt\pi}\bar m\,s^{3/2}+2\beta s\).
\end{lemma}
\begin{proof}
Let \(Z(y,r)=\mathbb E|y+\sqrt{2r}\,N|\), \(N\) standard normal. It solves \(Z_r=Z_{yy}\), with \(Z(y,0)=|y|\), \(Z_y(0,r)=0\),
\(Z_y\ge0\) on \(y\ge0\), and \(Z(0,r)=2\sqrt{r/\pi}\). Put \(\zeta(x,r)=Z(d(x),r)\). For \(\varphi\ge0\) in \(H^1(\Omega)\) we get
\(\avg{\zeta_r\varphi}+\avg{\nabla\zeta\cdot\nabla\varphi}=\avg{\nabla d\cdot\nabla(Z_y(d,r)\varphi)}\ge0\): the function
\(Z_y(d,r)\varphi\) is non-negative and vanishes on \(\partial\Omega\), and \(d\) is concave. So \(\zeta\) is a supersolution and
\(P_rd\le\zeta\), in particular \(P_rd\le2\sqrt{r/\pi}\) on \(\partial\Omega\). Then \(P_r\ut\le\bar m\,2\sqrt{r/\pi}+\beta\) on
\(\partial\Omega\). Insert this in \zcref{lem:layer}, using \(|\Omega|^{-1}\oint m=2\).
\end{proof}

\begin{proposition}\label{prop:G}
With \(E:=c^2-\avg{u^2}\), for all \(a,b\ge0\),
\[
G(a,b)\ge\avg{u^3}+2E(a+b)+4abc-(c+2b)\Lambda(a)-(c+2a)\Lambda(b)-c\Lambda(a+b).
\]
\end{proposition}
\begin{proof}
\(G(a,b)=\avg{\ut P_auP_bu}-cF(a+b)\). Write \(P_au=u-2a+\Gamma_a\) and expand. Since \(\ut\ge0\) and \(\Gamma\ge0\):
\(\avg{\ut\Gamma_a\Gamma_b}\ge0\); \(\avg{\ut(u-2a)\Gamma_b}=\avg{\ut^2\Gamma_b}-(c+2a)\Lambda(b)\ge-(c+2a)\Lambda(b)\), and likewise with
\(a,b\) exchanged; and \(\avg{\ut(u-2a)(u-2b)}=\avg{u^3}+c\avg{u^2}-2(a+b)\avg{u^2}+4abc\). Finally use
\(F(a+b)=\avg{u^2}-2c(a+b)+\Lambda(a+b)\).
\end{proof}

\begin{proposition}\label{prop:short}
Let \(\avg{u^3}\ge0\), \(\Lambda(s)\le Ls^{3/2}+Bs\), and let \(\tau_0>0\) satisfy
\begin{equation}\label{eq:S}
S(\tau_0):=\frac{E-cB}3-\frac8{35}cL\tau_0^{1/2}-\frac{128}{3465}L\tau_0^{3/2}-\frac B{15}\tau_0>0 .
\end{equation}
Then \(K(\tau)\ge\frac13\avg{u^3}\tau^3+S(\tau_0)\tau^4>0\) on \((0,\tau_0]\).
\end{proposition}
\begin{proof}
By~\eqref{eq:KG} and \zcref{prop:G}, integrate against \(2(\tau-a-b)\ge0\) on the triangle. Dirichlet integrals give the
weights \(2\iint(\tau-a-b)=\tau^3/3\), \(2\iint(\tau-a-b)(a+b)=\tau^4/6\), \(2\iint(\tau-a-b)4ab=\tau^5/15\),
\(2\iint(\tau-a-b)a^{3/2}=\frac{16}{315}\tau^{9/2}\), \(2\iint(\tau-a-b)(a+b)^{3/2}=\frac8{63}\tau^{9/2}\),
\(2\iint(\tau-a-b)ba^{3/2}=\frac{32}{3465}\tau^{11/2}\). Thus
\[
K\ge\tfrac13\avg{u^3}\tau^3+\tfrac13(E-cB)\tau^4+\tfrac c{15}\tau^5-\tfrac8{35}cL\tau^{9/2}-\tfrac{128}{3465}L\tau^{11/2}-\tfrac B{15}\tau^5,
\]
and \(S(\tau)\) decreases in \(\tau\).
\end{proof}

\begin{remark}\label{rem:E}
So \(K=\frac13\avg{u^3}\tau^3+\frac13E\tau^4+O(\tau^{9/2})\). For the triangle (inradius 1), \(E=9/100-27/700=9/175\), which
reproduces the coefficient in equation~(14) of the source, \(\mu_3\sim\frac{27}{19250}\mathrm{Pe}^3\tau^3+\frac9{175}\mathrm{Pe}^3\tau^4\). For
the disk \(E=1/24\), so \(K\sim\tau^4/72\). The point of \zcref{prop:short} is that \(E\) stays near \(1/24\) as \(n\to\infty\) while
\(\avg{u^3}\to0\): the window \((0,\tau_0]\) does not shrink.
\end{remark}

For a polynomial \(\ut_L\) with \(-\Delta\ut_L=2\) on \(\Pg\) let \(\theta_-\le\ut_L\le\theta_+\) on \(\partial\Pg\) and let \(\bar m\)
bound \(|\nabla\ut_L|\) on \(\partial\Pg\); then \(|\nabla\ut_L|\le\bar m\) on \(\Pg\), as \(|\nabla\ut_L|^2\) is subharmonic. The maximum
principle gives \(\ut\le\ut_L-\theta_-\) and \(\ut_L\le\theta_++\bar md\), so \zcref{lem:barrier} applies with
\(\beta=\theta_+-\theta_-\). Also \(c\) and \(\norm u\) are enclosed by the corresponding quantities of \(\ut_L\), within
\(\theta_+-\theta_-\).

\section{The third moment}\label{sec:u3}

\begin{proposition}\label{prop:u3}
\(\avg{u^3}>0\) for every \(n\ge3\).
\end{proposition}
For \(n=3\) the value is \(27/19250\) (\zcref{sec:triangle}). For \(n\ge4\) we use the following.

\begin{lemma}\label{lem:Hsmall}
On \(\Pg\) let \(Y=\frac12(1-r^2)+t^2/6\) and \(H=\ut-Y\). Then \(H\) is harmonic, \(H=\sigma^2/2-t^2/6\) on the side at arclength
\(\sigma\) from its midpoint, \(-t^2/6\le H\le t^2/3\), and \(H_c=H-\avg H\) satisfies
\(\norm{H_c}^2\le\frac{t^4}{16}\bigl(\frac1{n+1}+\frac43t^2\bigr)\).
\end{lemma}
\begin{proof}
On the unit circle, inside \(\Pg\), \(|H-t^2/12|\le t^2/4\) by the maximum principle. There \(H=\beta_0+\sum_{k\ge1}\beta_kr^{nk}\cos(nk\theta)\),
so \(\frac12\sum\beta_k^2\le t^4/16\), and the circular mean of \((H-\beta_0)^2\) at radius \(r\le1\) is at most
\(r^{2n}t^4/16\). Integrating, \(\int_{r<1}(H-\beta_0)^2\le\frac{\pi}{n+1}\frac{t^4}{16}\). On \(\Pg\setminus\{r<1\}\), of area
\(nt-\pi\le nt\cdot t^2/3\), \(|H-\beta_0|\le t^2/2\). Divide by \(|\Pg|=nt\ge\pi\), and note \(\norm{H_c}\le\norm{H-\beta_0}\).
\end{proof}

Let \(Y_c=Y-\avg Y=\frac14+\frac{t^2}{12}-\frac{r^2}2\), \(\phi=Y_c^2-\avg{Y_c^2}\), \(\Phi\) the radial polynomial with
\(-\Delta\Phi=\phi\), and \(\varrho_\Phi=\Phi-\text{(mean of }\Phi\text{ on the side)}\) on \(\partial\Pg\). Since \(u=Y_c+H_c\), Green's identity gives
\[
\avg{u^3}=M(t)+\frac3{|\Pg|}\oint\varrho_\Phi\,\partial_nH+3\avg{Y_cH_c^2}+\avg{H_c^3},
\]
\[
M(t)=\avg{Y_c^3}-\frac3{|\Pg|}\oint H\,\partial_n\Phi=\frac{t^4}{360}+\frac{19t^6}{7560},
\]
the last equality being an exact polynomial identity. By \zcref{lem:harmonic}, \(|\frac3{|\Pg|}\oint\varrho_\Phi\partial_nH|\le6(t+R)t^5(q(t)/3)^{1/2}\), where
\(\int_0^t\varrho_\Phi^2=t^9q(t)\), \(q(0)=1/4082400\). By \zcref{lem:Hsmall} and \(1/(n+1)\le t/\pi\),
\(|3\avg{Y_cH_c^2}+\avg{H_c^3}|\le(\frac34+\frac54t^2+\frac12t^2)\frac{t^4}{16}(\frac t\pi+\frac43t^2)\). In interval arithmetic the
resulting lower bound for \(\avg{u^3}/t^4\) is at least \(1.06\cdot10^{-3}\) on \((0,\tan(\pi/40)]\) (\texttt{u3\_uniform.py}), so \(\avg{u^3}>0\) for
\(n\ge40\). The same bounds give \(\avg{u^3}=t^4/360+O(t^5)\). The script \texttt{sg\_uniform.py} of~\autocite{Result14}, based on an explicit Fourier approximant of \(H\), gives a second proof for \(n\ge140\).

For \(4\le n\le150\), let \(\ut_N\) be a polynomial with \(-\Delta\ut_N=2\), \(E=\ut-\ut_N\), \(f=E-\avg E\), \(u_N=\ut_N-\avg{\ut_N}\), and
\(\Phi\) a polynomial with \(-\Delta\Phi=u_N^2-\avg{u_N^2}\). Then
\(\avg{u^3}=\avg{u_N^3}+3|\Pg|^{-1}\oint(\Phi\partial_nE+\ut_N\partial_n\Phi)+3\avg{u_Nf^2}+\avg{f^3}\), and \zcref{lem:harmonic} bounds the
last terms and the term containing \(\partial_nE\). The script \texttt{sg\_certify.py} of~\autocite{Result14} evaluates this in ball
arithmetic; we re-ran it for this paper, and all \(147\) cases pass.

\section{Certified computation: \texorpdfstring{\(4\le n\le199\)}{4 <= n <= 199}}\label{sec:finite}

Fix \(\tau_0=0.015\). For each \(n\) the script \texttt{finite\_n.py} checks~\eqref{eq:S} and \(K>0\) on \([\tau_0,\infty)\).

\paragraph{Exact integration.} Let \(p=x^2+y^2\), \(q=\operatorname{Re}(x+\mathrm iy)^n\). Every \(D_n\)-invariant polynomial is a polynomial in \(p,q\). For a
homogeneous function \(F_d\) of degree \(d\), cone coordinates on \(T\) give
\(\avg{F_d}=\frac2{d+2}\frac1t\int_0^tF_d(1,\sigma)\dd\sigma\). The integrand is a polynomial in \(\sigma\), integrated exactly by Gauss--Legendre
rules with ball enclosures of nodes and weights. Thus all moments \(\avg{p^aq^k}\) are enclosed, and products reduce to these
moments. Also \(\Delta(p^aq^k)=4a(a+nk)p^{a-1}q^k+n^2k(k-1)p^{a+n-1}q^{k-2}\), and on the side \(\partial_xp=2\),
\(\partial_xq=n\operatorname{Re}(1+\mathrm i\sigma)^{n-1}\). All arithmetic is in Arb~\autocite{Arb} through python-flint, at 3000 or 5000 bits.

\paragraph{Approximants.} With \(h_k=\operatorname{Re}((x+\mathrm iy)/R)^{nk}\) we take
\(\ut_H=\frac12(1-p)+\sum_{k=0}^{K_H}a_kh_k\) and
\(g_H=\frac{p^2}{32}+\beta p+\sum_{k=1}^{K_H}(\gamma_kp+b_k)h_k\), with \(\gamma_k=-a_k/(4(1+nk))\), \(\beta=-\frac14(\frac12-\avg{\ut_H}+a_0)\), so that
\(-\Delta g_H=\ut_H-\avg{\ut_H}\) exactly. The coefficients are dyadic numbers from floating-point least squares on the side.
\zcref{lem:cell} involves only one-dimensional integrals, so \(K_H=24\) to \(40\) is cheap. A second pair \((\ut_L,g_L)\) of the same form,
with \(K_L\le8\), enters all two-dimensional integrals; \(\norm{g_H-g_L}\) is computed from the Gram matrix of the functions
\(p\), \(ph_k\), \(h_k\) by the cone formula. So \(\omega=\norm{g-g_L}\) is enclosed, with \(\hat g=g_L-\avg{g_L}\). We take \(w=\ut_L-\avg{\ut_L}\)
in \zcref{lem:compare}; then \(o=\theta_+-\theta_-\), and \(U\) comes from a grid of \(T\) plus the Lipschitz bound.

\paragraph{Modes.} On the span of the scaled monomials \((p/R^2)^a(q/R^n)^k\) with (\(k=0\), \(a\le A_0\)) or
(\(1\le k\le K_q\), \(a\le A_1\)) we compute Gram and stiffness matrices in ball arithmetic, a Cholesky factor, and the Ritz pairs in
double precision. We keep the \(M\) pairs with \(\mu_k\le800\) and subtract means; \(\mu_k,d_k\) and the coefficients are exact binary
numbers. How they were found is irrelevant for rigour. Then \(\dualnorm{\ell_k}\le C_P\norm{R_k-\avg{R_k}}+\tau_{\mathrm{tr}}\bdnorm{\partial_n\psi_k}\) with
\(R_k=-\Delta\psi_k-\mu_k\psi_k\), and \(\rho\), \(L_1\), \(\Delta_0\) of \zcref{lem:modal} and the matrix \(W_{kl}\) of~\eqref{eq:Ktilde} are enclosed.

\paragraph{The inequality.} On the grid \(\tau_{j+1}=1.02\,\tau_j\) (for \(\tau_j<0.2\)), \(\tau_{j+1}=1.1\,\tau_j\) (beyond), from \(\tau_0\) to \(40\), the
script encloses \(\tilde K,\tilde K',\tilde K''\) at \(\tau_j\) and verifies, with the bound~\eqref{eq:Taylor} on \(\tilde K'''\) and \(\mathcal E\) increasing,
that the lower Taylor bound on \([\tau_j,\tau_{j+1}]\) exceeds \(\mathcal E(\tau_{j+1})\). For \(\tau\ge40\) it verifies
\((2W_{00}-e_1)40>C_1\), with \(e_1\tau\ge\mathcal E(\tau)\) the linear bound of \zcref{lem:compare}. In particular \(2W_{00}-e_1>0\) gives \(\avg{ug^2}>0\).

\begin{proposition}\label{prop:finite}
For every \(4\le n\le200\), \(K>0\) on \([0.015,\infty)\) and \(S(0.015)>0\).
\end{proposition}

\zcref{tab:finite} shows sample data. The smallest value of \((\text{lower bound}-\mathcal E)/\tilde K\) over all \(n\) and grid points is
attained at \(\tau_0\); the minimum over all \(n\) is \(0.29\), at \(n=21\). The smallest \(S(\tau_0)\) is \(1.3\cdot10^{-4}\), at \(n=4\). The table gives \(\tilde K(\tau_0)\) and \(\mathcal E\) at the next grid point.

\begin{table}
\centering
\footnotesize
\setlength{\tabcolsep}{4pt}
\begin{tabular}{rrrrrrrrr}
\toprule
\(n\) & \(M\) & \(\omega\) & \(\rho\) & \(L_1\) & \(\Delta_0\) & \(S(\tau_0)\) & \(\tilde K(\tau_0)\) & \(\mathcal E\)\\
\midrule
4 & 39 & \(3.7\cdot10^{-8}\) & \(3.3\cdot10^{-6}\) & \(3.6\cdot10^{-7}\) & \(3.7\cdot10^{-6}\) & \(1.33\cdot10^{-4}\) & \(8.20\cdot10^{-10}\) & \(1.13\cdot10^{-10}\)\\
5 & 29 & \(1.4\cdot10^{-7}\) & \(3.6\cdot10^{-6}\) & \(1.3\cdot10^{-8}\) & \(3.9\cdot10^{-6}\) & \(7.91\cdot10^{-4}\) & \(5.51\cdot10^{-10}\) & \(1.14\cdot10^{-10}\)\\
6 & 23 & \(2.8\cdot10^{-7}\) & \(4.3\cdot10^{-6}\) & \(2.2\cdot10^{-13}\) & \(4.9\cdot10^{-6}\) & \(1.11\cdot10^{-3}\) & \(4.49\cdot10^{-10}\) & \(1.33\cdot10^{-10}\)\\
8 & 18 & \(5.0\cdot10^{-7}\) & \(4.2\cdot10^{-6}\) & \(9.5\cdot10^{-7}\) & \(6.1\cdot10^{-6}\) & \(1.46\cdot10^{-3}\) & \(3.73\cdot10^{-10}\) & \(1.51\cdot10^{-10}\)\\
12 & 13 & \(6.7\cdot10^{-7}\) & \(4.0\cdot10^{-6}\) & \(4.8\cdot10^{-6}\) & \(1.0\cdot10^{-5}\) & \(1.91\cdot10^{-3}\) & \(3.34\cdot10^{-10}\) & \(2.11\cdot10^{-10}\)\\
16 & 11 & \(7.9\cdot10^{-7}\) & \(4.0\cdot10^{-6}\) & \(4.4\cdot10^{-6}\) & \(1.0\cdot10^{-5}\) & \(2.11\cdot10^{-3}\) & \(3.24\cdot10^{-10}\) & \(2.03\cdot10^{-10}\)\\
24 & 9 & \(7.3\cdot10^{-7}\) & \(3.9\cdot10^{-6}\) & \(5.1\cdot10^{-6}\) & \(1.0\cdot10^{-5}\) & \(2.40\cdot10^{-3}\) & \(3.18\cdot10^{-10}\) & \(2.05\cdot10^{-10}\)\\
32 & 8 & \(5.2\cdot10^{-7}\) & \(3.9\cdot10^{-6}\) & \(3.9\cdot10^{-6}\) & \(8.8\cdot10^{-6}\) & \(2.64\cdot10^{-3}\) & \(3.16\cdot10^{-10}\) & \(1.66\cdot10^{-10}\)\\
48 & 8 & \(4.7\cdot10^{-7}\) & \(3.9\cdot10^{-6}\) & \(2.8\cdot10^{-6}\) & \(7.6\cdot10^{-6}\) & \(2.81\cdot10^{-3}\) & \(3.14\cdot10^{-10}\) & \(1.39\cdot10^{-10}\)\\
64 & 8 & \(2.9\cdot10^{-7}\) & \(3.9\cdot10^{-6}\) & \(1.6\cdot10^{-6}\) & \(6.1\cdot10^{-6}\) & \(2.93\cdot10^{-3}\) & \(3.14\cdot10^{-10}\) & \(1.09\cdot10^{-10}\)\\
100 & 8 & \(1.3\cdot10^{-7}\) & \(3.9\cdot10^{-6}\) & \(6.5\cdot10^{-7}\) & \(4.8\cdot10^{-6}\) & \(3.06\cdot10^{-3}\) & \(3.13\cdot10^{-10}\) & \(8.45\cdot10^{-11}\)\\
150 & 8 & \(6.4\cdot10^{-8}\) & \(3.9\cdot10^{-6}\) & \(2.9\cdot10^{-7}\) & \(4.3\cdot10^{-6}\) & \(3.13\cdot10^{-3}\) & \(3.13\cdot10^{-10}\) & \(7.47\cdot10^{-11}\)\\
199 & 8 & \(3.8\cdot10^{-8}\) & \(3.9\cdot10^{-6}\) & \(1.7\cdot10^{-7}\) & \(4.1\cdot10^{-6}\) & \(3.16\cdot10^{-3}\) & \(3.13\cdot10^{-10}\) & \(7.14\cdot10^{-11}\)\\
\bottomrule
\end{tabular}

\caption{Certificates of \zcref{prop:finite}, sample of \(n\). All numbers are rounded to the digits shown, \(\tau_0=0.015\).
Full data: \texttt{certificates/finite\_4\_200.json}.}
\label{tab:finite}
\end{table}

\section{All \texorpdfstring{\(n\ge200\)}{n >= 200} and the disk}\label{sec:large}

Now rescale \(\Pg\) to area \(\pi\): call it \(\mathcal P\), with apothem \(a\), \(a^2=\arctan(t)/t\), and circumradius \(v=aR\). The alternating
series of \(\arctan\) gives \(1-t^2/3\le a^2\le1-t^2/3+t^4/5\) and \(v^2\le1+\frac23t^2\). Time scales by \(a^2\). Let \(j_k\) be positive zeros of \(J_1\),
\(\mu_k=j_k^2\), \(\psi_k(r)=J_0(j_kr)/J_0(j_k)\), and on the unit disk \(D\)
\[
\bar u=\tfrac14-\tfrac12r^2,\qquad g_0=\tfrac1{32}r^4-\tfrac1{16}r^2+\tfrac1{48},\qquad -\Delta g_0=\bar u,\quad g_0'(1)=0,\quad\avg{g_0}_D=0.
\]
\begin{lemma}[note=disk data]\label{lem:disk}
In \(L^2(D)\) with the average product, the \(\psi_k\) are orthonormal, \(\avg{\bar u\psi_k}_D=-2/\mu_k\), \(\avg{g_0\psi_k}_D=-2/\mu_k^2\),
\(\avg{g_0^2}_D=1/11520\), \(\avg{\bar ug_0^2}_D=1/184320\), and
\[
\avg{\bar u\psi_k^2}_D=\frac1{12},\qquad\avg{\bar u\psi_k\psi_l}_D=-\frac{2(\mu_k+\mu_l)}{(\mu_k-\mu_l)^2}\quad(k\ne l),
\]
\[
\avg{\bar ug_0\psi_k}_D=\frac{\mu_k^2+96\mu_k-1728}{48\mu_k^3}.
\]
\end{lemma}
\begin{proof}
For a radial polynomial \(f\), Green's identity with \(\psi_k(1)=1\), \(\psi_k'(1)=0\) gives \(\avg{f\psi_k}_D=\mu_k^{-1}(2f'(1)-\avg{\psi_k\Delta f}_D)\),
which terminates. For the products, let \(E=\int_0^1r^2\psi_l\psi_k'\dd r\), \(F=\int_0^1r^2\psi_k\psi_l'\dd r\) with \(k\ne l\). Integration by parts gives \(E+F=1\). Multiplying
the two Bessel equations by \(r^2\psi_l'\), \(r^2\psi_k'\) and adding gives \(\mu_kF+\mu_lE=0\). Also \((\mu_k-\mu_l)2\int_0^1r^3\psi_k\psi_l\dd r=4(E-F)\). So
\(\avg{r^2\psi_k\psi_l}_D=4(\mu_k+\mu_l)/(\mu_k-\mu_l)^2\). The Pohozaev identity \(\int_0^1r^4\psi_k'(\psi_k''+\psi_k'/r+\mu_k\psi_k)\dd r=0\) gives
\(\avg{r^2\psi_k^2}_D=1/3\).
\end{proof}

With \(D_0=1\), \(\psi_0=g_0\), \(\mu_0=0\), \(D_k=2/\mu_k^2\), \(M=12\), let \(W_{kl}=D_kD_l\avg{\bar u\psi_k\psi_l}_D\) and define \(K_D\) by~\eqref{eq:Ktilde}. This is a
fixed function, independent of \(n\). On \(\mathcal P\) we use \(\hat g=g_0-\avg{g_0}\) and \(\psi_k-\avg{\psi_k}\), with averages over \(\mathcal P\). Set \(x=t^2\).

\begin{lemma}\label{lem:largeerrors}
On \(\mathcal P\), let \(C_P=v/\pi\), \(\tau_{\mathrm{tr}}^2=(2C_P^2+2vC_P)/a\). The following bounds hold, every right-hand side being increasing in \(t\):
\begin{enumerate}
\item \(H=\ut-\frac12(1-r^2)\) is harmonic, its oscillation on \(\mathcal P\) is at most \(x/2\), and \(\norm{H_c}\le\frac x4(t/\pi+\frac43x)^{1/2}\);
\item the flux of \(g_0\) is \(e_0=\frac a8(r^2-1)\) on the side, and \(\bdnorm{e_0}^2\le x^2Q/32\), where \(Q=\frac{A^2}5+\frac23A(\zeta-\frac13)+(\zeta-\frac13)^2\), \(A=a^2=1-\frac x3+x\zeta\),
\(0\le\zeta\le x/5\);
\item \(\omega:=\norm{g-\hat g}\le C_P(C_P\norm{H_c}+\tau_{\mathrm{tr}}\bdnorm{e_0})\);
\item the flux of \(\psi_k\) satisfies \(|\partial_n\psi_k|\le16\,\Psi_k|e_0|\), where \(\Psi_k=\sup|\Psi''|\) over \([a^2,v^2]\), \(\Psi(r^2)=\psi_k(r)\); so \(\dualnorm{\ell_k}\le16\tau_{\mathrm{tr}}\Psi_k\bdnorm{e_0}\);
\item with \(\zeta_M=g_0-\sum_{k\le M}d_k\psi_k\), \(d_k=-2/\mu_k^2\): \(\rho^2\le\frac1{11520}-\sum_{k\le M}d_k^2+\frac x3\sup_{[1,v]}\zeta_M^2\);
\item for radial Lipschitz \(f\): \(|\avg f_{\mathcal P}-\avg f_D|\le w_1\operatorname{Lip}f\), \(w_1=x^2/6\); and \(\avg{\bar u}_{\mathcal P}=-\bar b/2\), \(0\le\bar b=\frac12(a^2(1+\frac x3)-1)\le x^2/20\).
\end{enumerate}
\end{lemma}
\begin{proof}
(1) As in \zcref{lem:Hsmall}, on \(\mathcal P\), with inscribed circle of radius \(a\). The wall values of \(H\) are \(\frac12(r^2-1)\). (2) \(g_0'=r(r^2-1)/8\), and on the
side \(\partial_x=(a/r)\partial_r\); then integrate. (3) \zcref{lem:cell}'s proof: here \(u-(\bar u-\avg{\bar u})=H_c\) and \(\avg{\nabla(g-\hat g)\cdot\nabla\varphi}=\avg{H_c\varphi}-|\mathcal P|^{-1}\oint e_0\varphi\).
(4) \(\partial_n\psi_k=2a\Psi'(r^2)\) and \(\Psi'(1)=0\), so \(|\partial_n\psi_k|\le2a\Psi_k|r^2-1|=16\Psi_k|e_0|\); \(\psi_k\) solves the Helmholtz equation exactly, so \(\ell_k\) is only the boundary term.
(5) \(\avg{\zeta_M^2}_D=\avg{g_0^2}_D-\sum d_k^2\), by orthonormality; and \(|\mathcal P\setminus D|\le\pi(1-a^2)\le\pi x/3\). (6) The laws of \(r\) under \(\mathcal P\) and \(D\) agree on
\([0,a]\) and their distribution functions differ by at most \(1-a^2\le x/3\) on \([a,v]\), with \(v-a\le x/2\).
\end{proof}

\begin{proposition}\label{prop:large}
For \(N=200\), every \(n\ge N\), and \(\hat\tau_0=0.0149\): \(K_n>0\) on \([\hat\tau_0,\infty)\) in the time unit of \(\mathcal P\), and \(S(0.015)>0\) in the apothem-\(1\) frame. The same holds for the disk.
\end{proposition}
\begin{proof}
By \zcref{lem:modal,lem:compare} with \(w=\bar u-\avg{\bar u}_{\mathcal P}\), \(U=\frac14+\frac x3+\frac{\bar b}2\), \(o=x/2\), and by \zcref{lem:largeerrors}(6), which replaces averages over \(\mathcal P\) by averages over \(D\),
\[
K_n(\tau)\ge K_D(\tau)-\mathcal E(\tau)-r_W\tau,\qquad r_W=2\bigl(w_1(vF^2+4UFL)+\tfrac12\bar bF^2\bigr),
\]
where \(F\ge\sup|f_s|\), \(L\ge\sup|f_s'|\) on \([0,v]\) for \(f_s=g_0-\sum_kd_k\mathrm e^{-\mu_ks}\psi_k\); we use \(|J_0|\le1\), \(|J_1|\le1/\sqrt2\). For the term \(r_W\), note
\(\avg{w\tilde A_a\tilde A_b}_{\mathcal P}\) differs from \(\avg{\bar uf_af_b}_D\) by at most \(w_1\operatorname{Lip}(\bar uf_af_b)+\frac{\bar b}2F^2+2w_1LUF\).
All constants are increasing in \(t\); \texttt{uniform.py} evaluates them at \(t=\tan(\pi/200)\): \(\omega\le4.2\cdot10^{-6}\), \(L_1\le1.6\cdot10^{-6}\), \(\rho\le1.9\cdot10^{-6}\), and checks the
inequality on the grid of \zcref{sec:finite} and for \(\tau\ge40\). At \(\hat\tau_0\): \(K_D=3.06\cdot10^{-10}\), error \(2.06\cdot10^{-10}\). For the disk, \(t=0\), all errors vanish except the truncation \(\rho\le1.1\cdot10^{-6}\).
In the apothem-\(1\) frame, for \(S\): \(H'=\ut-Y\) has \(-x/6\le H'\le x/3\), so \(c\in[\frac14-\frac x{12},\frac14+\frac{5x}{12}]\), \(\norm{\ut}\le\norm Y+x/3\), \(E=2c^2-\norm{\ut}^2\), and \(\ut\le Rd+x/2\) because \(Y\le x/6+Rd\); so
\(\bar m=R\), \(\beta=x/2\).
\end{proof}

\begin{proof}[of \zcref{thm:main}]
Let \(n\ge4\). \(\avg{u^3}>0\) by \zcref{prop:u3}. By \zcref{prop:short} with \zcref{prop:finite} (\(n\le199\)) or \zcref{prop:large} (\(n\ge200\)), \(K>0\) on \((0,0.015]\) in the apothem-\(1\) frame. On
\([0.015,\infty)\) use \zcref{prop:finite}; for \(n\ge200\), \zcref{prop:large} covers \(\tau\ge\hat\tau_0/a^2\), and \(\hat\tau_0/a^2\le0.015\) since \(a^2\ge1-t^2/3\). The disk is the same with \(\avg{u^3}=0\), \(\bar m=1\), \(\beta=0\), \(c=\frac14\), \(E=\frac1{24}\).
\end{proof}

\section{The equilateral triangle}\label{sec:triangle}

For \(n=3\), \(R=2\), and everything except the modes is an exact polynomial in \(p\) and \(q=\operatorname{Re}(x+\mathrm iy)^3\):
\(\ut=\frac12(1-p)+\frac16(1-q)\), \(c=\frac3{10}\), \(\avg{u^2}=\frac{27}{700}\), \(\avg{u^3}=\frac{27}{19250}\), \(-\frac3{10}\le u\le\frac{11}{30}\), and
\[
g=-\frac{p^3}{960}+\frac{p^2}{32}-\frac{11p}{120}+\frac{pq}{96}+\frac{q^2}{480}+\text{const},\qquad\avg{ug}=\frac9{1540},\quad\avg{ug^2}=-\frac{7101}{476476000}.
\]
So \(\omega=o=0\). The long-time coefficient in the normalisation of the source is \(3\avg{ug_2}/(2\avg{ug})^{3/2}=-789\sqrt{770}/309400=-0.0707623\dots\), and the geometric skewness is
\(\avg{u^3}/\avg{u^2}^{3/2}=4\sqrt{21}/99\).

For derivatives: with \(e=A-\tilde A\) one has
\(K'-\tilde K'=2\int_0^\tau\avg{ue(s)(P_{\tau-s}u+\tilde A'(\tau-s))}\dd s+2\avg{u\tilde A(0)e(\tau)}\), by differentiating the identity of \zcref{lem:compare}. So, by~\eqref{eq:PW},
\begin{equation}\label{eq:Kprime}
|K'-\tilde K'|\le2U\Delta_0\Bigl(\frac{\norm u}\nu+\sum_k|d_k|\norm{\psi_k}+\rho\Bigr).
\end{equation}
\begin{proof}[of \zcref{thm:triangle}]
\zcref{prop:short} with \(\bar m=\frac32\), \(\beta=0\), \(L=8\cdot\frac32/(3\sqrt\pi)\): the function
\(h(x)=\frac13\avg{u^3}+\frac E3x^2-\frac8{35}cLx^3-\frac{128}{3465}Lx^5\), \(x=\sqrt\tau\), increases and then decreases, and \(h(0),h(\sqrt{0.025})>0\); so \(K>0\) on \((0,0.025]\). With
\(35\) Ritz modes (\(\mu_k\le400\), \(\rho+L_1\le1.5\cdot10^{-5}\)), \texttt{triangle.py} verifies \(K>0\) on \([0.025,1]\), \(K'<0\) on \([1,1.6]\) by~\eqref{eq:Kprime}, \(K<0\) on \([1.6,\infty)\), and locates the zero.
\end{proof}

\section{Table 1 of the source}\label{sec:table}

\zcref{tab:source} compares the printed values (arXiv v3; circumradius 1) with enclosures from~\autocite{Result14}, script \texttt{table1.py}, re-run for this paper. All printed values agree up to rounding, except the arXiv-only
\(S^G_{100}=8.6\cdot10^{-7}\), about 3\% too large; it does not appear in the published table.

\begin{table}
\centering
\footnotesize
\setlength{\tabcolsep}{4pt}
\begin{tabular}{rllll}
\toprule
\(n\) & long-time, source & enclosure & \(S^G\), source & enclosure \\
\midrule
3 & \(-0.035\) & \(-0.0353812\dots\) (exact) & 0.185 & \(0.1851546\dots\) (exact)\\
4 & 0.139 & \([0.13849,0.13856]\) & 0.0813 & \([0.081294,0.081297]\)\\
5 & 0.194 & \([0.19390,0.19429]\) & 0.042 & \([0.042262,0.042270]\)\\
6 & 0.215 & \([0.21474,0.21603]\) & 0.024 & \([0.024264,0.024274]\)\\
7 & 0.225 & \([0.22384,0.22699]\) & 0.0149 & \([0.014939,0.014955]\)\\
8 & 0.231 & \([0.22903,0.23281]\) & 0.0097 & \([0.0097036,0.0097168]\)\\
20 & 0.243 & \([0.23937,0.24655]\) & \(3.97\cdot10^{-4}\) & \([3.9611,3.9832]\cdot10^{-4}\)\\
100 & 0.245 & \([0.24384,0.24589]\) & \(8.6\cdot10^{-7}\) & \([8.3475,8.3971]\cdot10^{-7}\)\\
\(\infty\) & 0.245 & \(\sqrt6/10=0.2449490\) & 0 & 0\\
\bottomrule
\end{tabular}
\caption{Long-time coefficient \(3\avg{ug_2}/(2\avg{ug})^{3/2}\) and geometric skewness \(S^G=\avg{u^3}/\avg{u^2}^{3/2}\), circumradius 1.}
\label{tab:source}
\end{table}

\section{Computer verification and limitations}\label{sec:computer}

All scripts are in \texttt{all-times/scripts/} unless noted, print PASS or FAIL per claim, and exit non-zero on failure.
Environment: Python 3.12, python-flint 0.9.0 (Arb), numpy, scipy, sympy, under WSL Fedora. \texttt{run\_all.sh} runs everything; logs are in \texttt{all-times/logs/}.

\begin{center}
\footnotesize
\begin{tabular}{lp{5.2cm}l}
\toprule
script & claim & time \\
\midrule
\texttt{run\_finite.py 4 200} & \zcref{prop:finite} (\texttt{finite\_n.py}, \texttt{modal\_lib.py}) & 45 min, 3 cores \\
\texttt{uniform.py 200}, \texttt{uniform.py 0} & \zcref{prop:large}, large \(n\), the disk & 10 s \\
\texttt{u3\_uniform.py 40} & \(\avg{u^3}>0\), \(n\ge40\) & 5 s \\
\texttt{../scripts/sg\_certify.py 4 150} & \(\avg{u^3}>0\), \(4\le n\le150\) & 15 min \\
\texttt{../scripts/sg\_uniform.py 140 6} & \(\avg{u^3}>0\), \(n\ge140\), second proof & 5 s \\
\texttt{triangle.py} & \zcref{thm:triangle} & 2 min \\
\texttt{../scripts/triangle\_exact.py} & exact triangle values & 5 s \\
\texttt{../scripts/table1.py} & \zcref{tab:source} & 5 min \\
\texttt{numerics\_check.py} & cross-checks, not a certificate & 3 min \\
\bottomrule
\end{tabular}
\end{center}

\paragraph{Cross-checks.} For the square, let \(\hat u_\xi\) be the Fourier coefficients of \(u\) on the torus \((\mathbb R/2\mathbb Z)^2\). Then
\[
G(a,b)=\sum_{\xi,\eta\in\mathbb Z^2}\hat u_\xi\hat u_\eta\hat u_{\xi+\eta}\mathrm e^{-\pi^2(|\xi|^2a+|\eta|^2b)},
\]
and the resulting \(K\) agrees with the Galerkin value to four digits at \(\tau=0.015\) (\(8.20\cdot10^{-10}\)), and with the certified \(\tilde K\). The law in \zcref{rem:E} is confirmed numerically and against equation~(14) of the source.

\paragraph{Trusted base.} Payne--Weinberger constant~\autocite{PayneWeinberger1960,Bebendorf2003}; \(H^2\) regularity on convex polygons~\autocite{Grisvard1985}; the spectral theorem for the Neumann Laplacian; Green's identities;
Arb and python-flint for ball arithmetic; sympy for exact rational identities. Floating point only chooses coefficients.

\paragraph{Limitations.} The proof is computer-assisted and has not been verified independently. The window \(\tau_0=0.015\) and the threshold \(N=200\) are technical. Numerics suggest more than we prove:
\(G(a,b)>0\) on the whole quadrant for every \(n\ge4\), so that \(K\) is convex (and \(\avg{ue_1^2}>0\) for the first invariant eigenfunction, already in the square, where it is about \(+0.023\)); we did not prove this.
The monotonicity in \(n\) suggested by Table~1 of the source is not addressed, nor are the unimodality and median statements of the source.

\paperreferences
\end{document}
